\documentclass[10pt,a4paper]{amsart}
\usepackage[margin=19mm]{geometry}
\usepackage{amsmath,amssymb,mathtools}
\usepackage[scr=rsfs,cal=euler]{mathalfa}
\usepackage{xfrac}
\usepackage[unicode,hidelinks,bookmarks=false]{hyperref}
\newcommand{\ie}{i.e.\ }
\newcommand{\cf}{cf.\ }
\newcommand{\resp}{respectively\ }
\newcommand{\Subsection}{\subsection}
\providecommand{\nobreakdash}{\nobreak\hbox{-}}
\setlength{\emergencystretch}{2em}
\multlinegap0pt
\usepackage{amssymb}
\newtheorem{theorem}{Theorem}
\newcommand{\R}{\mathbb{R}}
\title{Convexity of the embedding parameter sets of some analytic function spaces}
\author{Benoit F. Sehba}
\date{Source: arXiv:2605.06277v1 (CC BY 4.0)}
\begin{document}
\maketitle

\noindent\emph{Source and licence.} Convexity of the embedding parameter sets of some analytic function spaces. Version arXiv:2605.06277v1 (CC BY 4.0), distributed by arXiv under the Creative Commons Attribution 4.0 International licence, http://creativecommons.org/licenses/by/4.0/, as recorded in the arXiv licence metadata for this exact version and shown on the abstract page https://arxiv.org/abs/2605.06277v1. Full licence notice: \texttt{licenses/arxiv-2605.06277-CC-BY-4.0.txt}.\par\smallskip

\noindent\textit{Editorial scope.} Counted unit: the article's theorem thm:convexity (source0.tex lines 339--352). The article's complete original proof of it is delivered with it (source0.tex lines 354--444, one \texttt{proof} environment of 91 lines). No mathematics is written, repaired or re-derived: the statement and the proof are the article's own lines, in the article's own notation, and no retained line is substituted or re-flowed.  No further source-local context is needed: every cross-reference of every retained line resolves inside the two delivered spans.  Nothing was omitted from any retained span, so the package carries no omission record.  No retained line contains a citation, so the wrapper carries no bibliography and no bibliography entry.  No retained line contains a figure environment, an image-inclusion command or drawing code, so no image is delivered and the package declares no asset dependency.  Editorial formatting only, in addition: an AMS \texttt{amsart} preamble that restates the article's own declarations and macros used by the retained lines, namely the environments \texttt{theorem} on separate counters, together with the standard packages those lines need.  The retained environments print in this document as Theorem 1 in order of appearance, whereas the article prints the same items with the numbering of its own full document.  The complete native source of the article as distributed in the arXiv e-print for this exact version is bundled unchanged as \texttt{sources/200006/source0.tex}.
\begin{theorem}\label{thm:convexity}
Let $\Phi_1$ be a growth function such that $\Phi_1^{-1}\in \Delta_{\log}^-$ with constant $C_1\geq 1$, and let $\Phi_2$ be a growth function such that $\Phi_2^{-1}\in \Delta_{\log}^+$ with constant
$C_2 \ge 1$. Then
the set
\[
\mathcal{E} = \left\{ (\alpha,\beta) \in \R^2 \;\middle|\;
\alpha, \beta \ge -1, \quad
\exists\, C,K > 0 \;\text{ s.t. }\;
\Phi_1^{-1}(t^{2+\alpha}) \le C\,\Phi_2^{-1}(Kt^{2+\beta})
\;\; \forall\, t > 0
\right\}
\]
is convex.
\end{theorem}

\begin{proof}
Let $(\alpha_0, \beta_0), (\alpha_1, \beta_1) \in \mathcal{E}$. By definition,
there exist constants $M_0, M_1 > 0$ and $K_0, K_1 > 0$ such that for all $t > 0$,
\begin{align}
\Phi_1^{-1}\!\left( t^{2+\alpha_0} \right) &\le M_0\,\Phi_2^{-1}\!\left( K_0 t^{2+\beta_0} \right), \label{eq:e0}\\
\Phi_1^{-1}\!\left( t^{2+\alpha_1} \right) &\le M_1\,\Phi_2^{-1}\!\left( K_1 t^{2+\beta_1} \right). \label{eq:e1}
\end{align}
Fix $\theta \in [0,1]$ and set
\[
\alpha_\theta = (1-\theta)\alpha_0 + \theta\alpha_1, \qquad
\beta_\theta = (1-\theta)\beta_0 + \theta\beta_1.
\]
Let us show that $(\alpha_\theta, \beta_\theta) \in \mathcal{E}$.

For any $t > 0$, write
\[
t^{2+\alpha_\theta}
= \left(t^{2+\alpha_0}\right)^{1-\theta}
\cdot \left(t^{2+\alpha_1}\right)^{\theta}.
\]
Since $\Phi_1^{-1}\in\Delta_{\log}^-$ with constant $C_1\ge 1$, it holds that for all $u,v>0$,
\[
\Phi_1^{-1}\!\left(u^{1-\theta}v^{\theta}\right) \le C_1^{\theta(1-\theta)}\,
\Phi_1^{-1}(u)^{1-\theta}\,\Phi_1^{-1}(v)^{\theta}.
\]
Applying this with $u = t^{2+\alpha_0}$ and $v = t^{2+\alpha_1}$ gives
\begin{equation}\label{eq:step1}
\Phi_1^{-1}(t^{2+\alpha_\theta})
\le C_1^{\theta(1-\theta)}\,
\Phi_1^{-1}(t^{2+\alpha_0})^{1-\theta}\,
\Phi_1^{-1}(t^{2+\alpha_1})^{\theta}.
\end{equation}

Raising \eqref{eq:e0} to the power $(1-\theta)$ and \eqref{eq:e1}
to the power $\theta$, gives 
\begin{align}
\Phi_1^{-1}(t^{2+\alpha_0})^{1-\theta}
&\le M_0^{1-\theta}\,\Phi_2^{-1}(K_0 t^{2+\beta_0})^{1-\theta},\label{eq:pow0}\\
\Phi_1^{-1}(t^{2+\alpha_1})^{\theta}
&\le M_1^{\theta}\,\Phi_2^{-1}(K_1 t^{2+\beta_1})^{\theta}.\label{eq:pow1}
\end{align}
Then multiplying \eqref{eq:pow0} and \eqref{eq:pow1}, we get
\begin{equation}\label{eq:step2}
\Phi_1^{-1}(t^{2+\alpha_0})^{1-\theta}\,
\Phi_1^{-1}(t^{2+\alpha_1})^{\theta}
\le M_0^{1-\theta} M_1^{\theta}\,
\Phi_2^{-1}(K_0 t^{2+\beta_0})^{1-\theta}\,
\Phi_2^{-1}(K_1 t^{2+\beta_1})^{\theta}.
\end{equation}

Since $\Phi_2^{-1}\in\Delta_{\log}^+$ with constant $C_2\ge 1$, it holds that for all $u,v>0$,
\[
\Phi_2^{-1}(u)^{1-\theta}\Phi_2^{-1}(v)^{\theta}
\le 
\Phi_2^{-1}\!\left(C_2^{\theta(1-\theta)}\,u^{1-\theta}v^{\theta}\right).
\]
We apply this with $u = K_0 t^{2+\beta_0}$ and $v = K_1 t^{2+\beta_1}$. For this, we observe that
\[
u^{1-\theta}v^{\theta}
= (K_0 t^{2+\beta_0})^{1-\theta} (K_1 t^{2+\beta_1})^{\theta}
= K_0^{1-\theta}K_1^{\theta} \,
t^{(1-\theta)(2+\beta_0)+\theta(2+\beta_1)},
\]
and $(1-\theta)(2+\beta_0)+\theta(2+\beta_1) = 2+\beta_\theta$. Therefore
\[
u^{1-\theta}v^{\theta} = L_\theta \, t^{2+\beta_\theta},
\]
where we set $L_\theta := K_0^{1-\theta}K_1^{\theta}$. Consequently,
\begin{equation}\label{eq:step3}
\Phi_2^{-1}(K_0 t^{2+\beta_0})^{1-\theta}\,
\Phi_2^{-1}(K_1 t^{2+\beta_1})^{\theta}
\le 
\Phi_2^{-1}\!\left(C_2^{\theta(1-\theta)}\,L_\theta t^{2+\beta_\theta}\right).
\end{equation}

Substituting \eqref{eq:step3} into \eqref{eq:step2} yields
\[
\Phi_1^{-1}(t^{2+\alpha_0})^{1-\theta}\,
\Phi_1^{-1}(t^{2+\alpha_1})^{\theta}
\le M_0^{1-\theta} M_1^{\theta} C_2^{\theta(1-\theta)}\,
\Phi_2^{-1}\!\left(C_2^{\theta(1-\theta)}\,L_\theta t^{2+\beta_\theta}\right).
\]
Finally, inserting this into \eqref{eq:step1} gives
\[
\Phi_1^{-1}(t^{2+\alpha_\theta})
\le C_1^{\theta(1-\theta)} M_0^{1-\theta} M_1^{\theta} C_2^{\theta(1-\theta)}\,
\Phi_2^{-1}\!\left(C_2^{\theta(1-\theta)}\,L_\theta t^{2+\beta_\theta}\right)
= C_\theta \, \Phi_2^{-1}\!\left(K_\theta t^{2+\beta_\theta}\right),
\]
where $C_\theta := M_0^{1-\theta} M_1^{\theta} C_1^{\theta(1-\theta)}$ and $K_\theta=C_2^{\theta(1-\theta)}\,L_\theta$. This shows $(\alpha_\theta,\beta_\theta) \in \mathcal{E}$. The proof is complete.
\end{proof}

\end{document}
